\documentclass[10pt,a4paper]{amsart}
\usepackage[margin=19mm]{geometry}
\usepackage{amsmath,amssymb,mathtools}
\usepackage[scr=rsfs,cal=euler]{mathalfa}
\usepackage{xfrac}
\usepackage[unicode,hidelinks,bookmarks=false]{hyperref}
\newcommand{\ie}{i.e.\ }
\newcommand{\cf}{cf.\ }
\newcommand{\resp}{respectively\ }
\newcommand{\Subsection}{\subsection}
\providecommand{\nobreakdash}{\nobreak\hbox{-}}
\setlength{\emergencystretch}{2em}
\multlinegap0pt
\usepackage[shortlabels]{enumitem}
\usepackage{csquotes}
\usepackage{tcolorbox}
\usepackage{tikz}
\usepackage{amssymb}
\usepackage{mathtools}
\usepackage{xcolor}
\newtheorem{lemma}{Lemma}
\newtheorem{proposition}{Proposition}
\newtheorem{theorem}{Theorem}
\newcommand{\Def}{\mathrm{Def}}
\newcommand{\R}{\mathbb{R}}
\newcommand{\Scal}{\mathrm{Scal}}
\title{Sharp thresholds for the Escobar functional: the Escobar--Willmore mass, geometric selection, and compactness trichotomy}
\author{Mayukh Mukherjee, Utsab Sarkar}
\date{Source: arXiv:2601.22665v2 (CC BY 4.0)}
\begin{document}
\maketitle

\noindent\emph{Source and licence.} Sharp thresholds for the Escobar functional: the Escobar--Willmore mass, geometric selection, and compactness trichotomy. Version arXiv:2601.22665v2 (CC BY 4.0), distributed by arXiv under the Creative Commons Attribution 4.0 International licence, http://creativecommons.org/licenses/by/4.0/, as recorded in the arXiv licence metadata for this exact version and shown on the abstract page https://arxiv.org/abs/2601.22665v2. Full licence notice: \texttt{licenses/arxiv-2601.22665-CC-BY-4.0.txt}.\par\smallskip

\noindent\textit{Editorial scope.} Counted unit: the article's theorem thm:intrinsic-single-gap (source0.tex lines 2070--2093). The article's complete original proof of it is delivered with it (source0.tex lines 2095--2263, one \texttt{proof} environment of 169 lines). No mathematics is written, repaired or re-derived: the statement and the proof are the article's own lines, in the article's own notation, and no retained line is substituted or re-flowed.  Retained uncounted as the source-local context the proof needs: lines 991--996; lines 1246--1248; lines 1250--1269; lines 1304--1315; lines 1772--1785; lines 1845--1859; lines 1994--2017; lines 2021--2026.  Nothing was omitted from any retained span, so the package carries no omission record.  No retained line contains a citation, so the wrapper carries no bibliography and no bibliography entry.  No retained line contains a figure environment, an image-inclusion command or drawing code, so no image is delivered and the package declares no asset dependency.  Editorial formatting only, in addition: an AMS \texttt{amsart} preamble that restates the article's own declarations and macros used by the retained lines, namely the environments \texttt{lemma}, \texttt{proposition}, \texttt{theorem} on separate counters, together with the standard packages those lines need.  The retained environments print in this document as Lemma 1, Proposition 1, Theorem 1 in order of appearance, whereas the article prints the same items with the numbering of its own full document.  The complete native source of the article as distributed in the arXiv e-print for this exact version is bundled unchanged as \texttt{sources/193658/source0.tex}.
\begin{lemma}[Global trace under a uniform collar]\label{lem:global-trace}
Assume \textup{(BG$^{2+}$)}. Then there exists $C_{\mathrm{tr}}<\infty$, depending only on the \textup{(BG$^{2+}$)} data, such that for all $u\in H^1(M)$,
\[
\|u\|_{L^2(\partial M)} + \|u\|_{L^{2^*_\partial}(\partial M)} \ \le\ C_{\mathrm{tr}}\|u\|_{H^1(M)}.
\]
\end{lemma}

\begin{lemma}[Uniform Fermi atlas and Jacobian control]\label{lem:atlas}
Assume \textup{(BG$^{2+}$)}. Then there exists an atlas of interior normal and boundary Fermi charts with uniform $C^2$ bounds and bilipschitz constants (independent of the center and chart), so that the Jacobians and their first derivatives are uniformly controlled on fixed chart radii. In particular, Euclidean capacity and trace estimates transfer to $(M,g)$ with constants depending only on the \textup{(BG$^{2+}$)} data (cf.\ Lemma~\ref{lem:global-trace}).
\end{lemma}

\begin{lemma}[Near-isometry of the energy with first-order geometric correction]\label{lem:near-iso-energy}
Let $(M^n,g)$ satisfy \textup{(BG$^{2+}$)} near a boundary point $x\in\partial M$.  Fix a Fermi chart $\Phi_x:B_{r_0}^+\to M$ adapted to the \emph{inner} unit normal, and for $0<\varepsilon<\varepsilon_0(R)$ use its restriction to $\varepsilon B_{2R}^+$.
 For $w\in H^1(\R^n_+)$ compactly supported in $B^+_{2R}$ and $u_\varepsilon:=\mathcal T_{x,\varepsilon}w$, we have
\begin{align}
\mathcal Q_g(u_\varepsilon)
&= a_n\int_{\R_+^n} |\nabla w|^2dy
+ a_n\varepsilon\int_{\R_+^n} y_n\Big(2\mathrm{II}^{ab}(x)\partial_{y_a}w\partial_{y_b}w - K_g(x)|\nabla w|^2\Big)dy \nonumber\\
&\quad + 2\varepsilon K_g(x)\|w(\cdot,0)\|_{L^2(\R^{n-1})}^2
+ O_R(\varepsilon^2)\|w\|_{H^1(B^+_{2R})}^2. \label{eq:Q-first-order}
\end{align}
Here $H_g=\frac{1}{n-1}\mathrm{tr}_{\bar g}\mathrm{II}$ is the mean curvature with respect to the inner normal. If, in addition, $w$ is tangentially radial, then the traceless part of $\mathrm{II}$ averages out and
\begin{align}
\mathcal Q_g(u_\varepsilon)
&= a_n\int_{\R_+^n} |\nabla w|^2dy \nonumber\\
&\quad + a_n\varepsilon K_g(x)\left[\frac{2}{n-1}\int_{\R_+^n} y_n|\nabla_{\tan} w|^2dy - \int_{\R_+^n} y_n|\nabla w|^2dy\right] \nonumber\\
&\quad + 2\varepsilon K_g(x)\|w(\cdot,0)\|_{L^2(\R^{n-1})}^2
+ O_R(\varepsilon^2)\|w\|_{H^1(B^+_{2R})}^2.
\label{eq:Q-isotropic}
\end{align}
\end{lemma}

\begin{lemma}[Boundary trace at critical scaling has no $O(\varepsilon)$ correction]\label{lem:trace-no-linear}
Let $n\ge 3$ and let $w\in H^1(\R^n_+)$ be supported in $B_{2R}^+$ with
$w(\cdot,0)\in L^{2^*_{\partial}}(\R^{n-1})$, where $2^*_{\partial}=\frac{2(n-1)}{n-2}$.
Assume \textup{(BG$^{2+}$)}. Then, for each fixed $R<\infty$ there exists
$\varepsilon_0(R)>0$ such that for all $0<\varepsilon<\varepsilon_0(R)$,
\[
\big\|\big(\mathcal T_{x,\varepsilon}w\big)\big|_{\partial M}\big\|_{L^{2^*_{\partial}}(\partial M)}^{2^*_{\partial}}
\ =\ \big(1+O_R(\varepsilon^2)\big)\|w(\cdot,0)\|_{L^{2^*_{\partial}}(\R^{n-1})}^{2^*_{\partial}}
\]
uniformly in $x\in\partial M$, where $O_R(\varepsilon^2)$ denotes a quantity bounded by
$C(R)\varepsilon^2$ with $C(R)$ depending only on $R$ and the \textup{(BG$^{2+}$)} data.
\end{lemma}

\begin{lemma}[Constrained second variation at $F_\ast$]\label{lem:hessian-sphere}
Let $q=\frac{2(n-1)}{n-2}$ and let $F_\ast$ be the Cayley pullback of a half-space optimizer, normalized by $\|F_\ast\|_{L^q(\partial S^n_+)}=1$. Then
\[
L_{g_{\mathrm{round}}}^\circ F_\ast=0\ \text{ in }S^n_+,
\qquad
B_{g_{\mathrm{round}}}^\circ F_\ast=\lambda_\ast F_\ast^{q-1}\ \text{ on }\partial S^n_+,
\]
with $\lambda_\ast=\mathcal N^\circ_{g_{\mathrm{round}}}(F_\ast)=\mathcal J_{g_{\mathrm{round}}}[F_\ast]=S_\ast$. For $\phi\in T_{F_\ast}\mathcal S=\{\phi:\int_{\partial S^n_+}F_\ast^{q-1}\phi=0\}$, the constrained second variation is
\begin{equation}\label{eq:Hessian-sphere}
Q[\phi]
= \mathcal N^\circ_{g_{\mathrm{round}}}(\phi,\phi)-(q-1)\lambda_\ast\int_{\partial S^n_+} F_\ast^{q-2}\phi^2d\sigma.
\end{equation}
On $(S^n_+,g_{\mathrm{round}})$ we have $H_{g_{\mathrm{round}}}\equiv0$, so the boundary mean-curvature term vanishes in $\mathcal N^\circ_{g_{\mathrm{round}}}$. Moreover, $\ker Q=E$ and there exists $c_\ast>0$ such that $Q[\phi]\ge c_\ast\|\phi\|_{H^1(S^n_+)}^2$ for all $\phi\in T_{F_\ast}\mathcal S\cap E^\perp$.
\end{lemma}

\begin{proposition}[Weighted Robin (Jacobi) spectral gap on $S^n_+$]\label{prop:robin-gap}
With $Q$ as in \eqref{eq:Hessian-sphere}, there exists $\gamma_\ast=\gamma_\ast(n)>0$ such that
\[
  Q[\phi]\ \ge\ \gamma_\ast\|\phi\|_{H^1(S^n_+)}^{2}
  \qquad\text{for all }\ \phi\in T_{F_\ast}\mathcal S\ \text{ with }\ \phi\perp E
  \text{ in the pairing } \langle\cdot,\cdot\rangle_{E}:=\mathcal N^\circ_{g_{\mathrm{round}}}(\cdot,\cdot).
\]
Equivalently, the Jacobi system
\[
  L_{g_{\mathrm{round}}}^\circ\phi=0\ \text{ in } S^n_+,
  \qquad
  B_{g_{\mathrm{round}}}^\circ\phi=(q-1)\lambda_\ast F_\ast^{q-2}\phi\ \text{ on } \partial S^n_+,
\]
has strictly positive first eigenvalue on $T_{F_\ast}\mathcal S\cap E^\perp$ (orthogonality in $\langle\cdot,\cdot\rangle_{E}$). By conformal invariance, the gap is uniform along $\mathcal M_S$.
\end{proposition}

\begin{lemma}[Graph norm equivalence and conformal transfer]\label{lem:graph-norm}
Let $\|\cdot\|_{E,g}$ denote the Escobar graph norm built from $(L_g^{\circ},B_g^{\circ})$ via the weak bilinear form
\[
\|u\|_{E,g}^2:=\mathcal N^\circ_g(u,u).
\]
\begin{enumerate}[label=(\alph*)]
\item (\emph{Equivalence on the sphere}) There exist $0<c\le C<\infty$ such that, for all $w\in H^1(S^n_+)$,
\[
c\|w\|_{H^1(S^n_+)}^2\ \le\ \|w\|_{E,g_{\mathrm{round}}}^2\ \le\ C\|w\|_{H^1(S^n_+)}^2.
\]
\item (\emph{Conformal transfer}) If $\widehat g=\phi^{\frac{4}{n-2}}g$ and $u=\phi v$ (equivalently $v=\phi^{-1}u$), then $\|v\|_{E,\widehat g}=\|u\|_{E,g}$ and $\Def^{\mathrm{Esc}}_{\widehat g}(v)=\Def^{\mathrm{Esc}}_{g}(u)$.
 In particular, for the Cayley transform $\mathcal C:S^n_+\to\R^n_+$ with its conformal factor,
\[
\|F-V^\sharp\|_{E,g_{\mathrm{round}}} = \|f-v^\sharp\|_{E,\mathrm{flat}},
\qquad \Def^{\mathrm{Esc}}_{g_{\mathrm{round}}}(F)=\Def^{\mathrm{Esc}}_{\mathrm{flat}}(f).
\]
\item (\emph{Explicit form and equivalence}) By the Green identity \eqref{eq:green-identity}, the graph norm is explicitly
\begin{equation}\label{eq:graph-norm-explicit}
\|u\|_{E,g}^2 = \int_M |\nabla u|^2dV_g + \frac{n-2}{4(n-1)}\int_M \Scal_gu^2dV_g + \frac{n-2}{2}\int_{\partial M} H_gu^2d\sigma_g.
\end{equation}
In particular, on $(S^n_+,g_{\mathrm{round}})$ where $H_{g_{\mathrm{round}}}\equiv 0$, this simplifies further.
Consequently, under \textup{(BG$^{2+}$)} and the coercivity assumption \textup{(Y$^+_{\partial}$)}, $\|\cdot\|_{E,g}$ and $\|\cdot\|_{H^1(M)}$ are equivalent.
\end{enumerate}
\end{lemma}

\begin{equation}\label{eq:green-identity}
\int_M uL_g^{\circ} vdV_g+\int_{\partial M} uB_g^{\circ} vd\sigma_g
= \int_M \langle \nabla u,\nabla v\rangle dV_g
 + \tfrac{n-2}{4(n-1)}\int_M \Scal_gu vdV_g
 + \tfrac{n-2}{2}\int_{\partial M} H_gu vd\sigma_g.
\end{equation}

\begin{theorem}[Intrinsic single-bubble coercivity on $(M,g)$, $n\ge5$]\label{thm:intrinsic-single-gap}
Assume $n\ge5$, \textup{(Y$^+_{\partial}$)} and \textup{(BG$^{2+}$)}. Equip $H^1(M)$ with the covariant weak energy pairing
\[
\langle u,v\rangle_{E,g}
:=\int_M\Big(\nabla u\cdot\nabla v+\tfrac{1}{a_n}\Scal_guv\Big)dV_g
  +\tfrac{n-2}{2}\int_{\partial M}H_guvd\sigma_g,
\qquad
\|u\|_{E,g}^2:=\langle u,u\rangle_{E,g}.
\]
Then there exist $\varepsilon_0>0$ and $c_\ast>0$ (depending only on $n$ and the \textup{(BG$^{2+}$)}, \textup{(Y$^+_{\partial}$)} data of $(M,g)$) such that for every $x\in\partial M$, every $\varepsilon\in(0,\varepsilon_0]$, and every $w\in H^1(M)$ obeying the \emph{constrained orthogonality}
\begin{equation}\label{eq:esc-orthog-gauge-final}
\langle w,\partial_\varepsilon U_{+,x,\varepsilon}\rangle_{E,g}
=\langle w,\nabla_x U_{+,x,\varepsilon}\rangle_{E,g}=0,
\qquad
\int_{\partial M} (U_{+,x,\varepsilon})^{q-1}wd\sigma_g=0,
\quad q=\tfrac{2(n-1)}{n-2},
\end{equation}
(one orthogonality condition in the weak pairing for each modulation direction and one boundary constraint fixing the amplitude) one has
\begin{equation}\label{eq:intrinsic-gap-final}
\big\langle D^2\mathcal J_g[U_{+,x,\varepsilon}]w,\ w\big\rangle
\ \ge\ c_\ast\|w\|_{E,g}^2
\ \simeq\ \|w\|_{H^1(M)}^2 .
\end{equation}
\end{theorem}

\begin{proof}[Proof (via blow-up and variational convergence)]
We argue by contradiction. Suppose there exist $x_k\in\partial M$, $\varepsilon_k\downarrow0$, and $w_k\in H^1(M)$ with
\begin{equation}\label{eq:contr-setup}
\|w_k\|_{E,g}=1,\qquad
\eqref{eq:esc-orthog-gauge-final}\ \text{ holds with }(x,\varepsilon)=(x_k,\varepsilon_k),
\qquad
\big\langle D^2\mathcal J_g[U_{+,x_k,\varepsilon_k}]w_k,\ w_k\big\rangle \longrightarrow 0 .
\end{equation}

\smallskip\noindent\textbf{Step 1 (Rescaling to $\R^n_+$ and nonvanishing of the weak limit).}
For each $k$, define
\[
\widetilde w_k(y)
:=\varepsilon_k^{(n-2)/2}w_k\big(\Phi_{x_k}(\varepsilon_k y)\big),
\qquad y\in B_{r_0/\varepsilon_k}^+\subset\R^n_+.
\]
By the near-isometry estimates (Lemma~\ref{lem:near-iso-energy}),
for each fixed $R\ge1$,
\begin{equation}\label{eq:per-ball-bound}
\|\nabla\widetilde w_k\|_{L^2(B_R^+)}^2
\le \|w_k\|_{E,g}^2+O_R(\varepsilon_k)=1+o_k(1),
\end{equation}
since the Fermi metric correction on $B_R^+$ (where $|y|\le R$)
is $O(\varepsilon_k R)$
and the scalar-curvature and mean-curvature contributions
in the chart are $O(\varepsilon_k^2 R^2)+O(\varepsilon_k)$
(both $\to0$ for $\varepsilon_k\to0$ with $R$ fixed).

We also need a \emph{global} gradient bound (on the full Fermi domain, not just fixed balls).
By the exact change of variables $z=\varepsilon_k y$ (which maps $B_{r_0/\varepsilon_k}^+$ to
$B_{r_0}^+$ in the Fermi chart),
\begin{equation}\label{eq:global-grad-bound}
\int_{B_{r_0/\varepsilon_k}^+}|\nabla_y\widetilde w_k|^2dy
=\int_{B_{r_0}^+(z)}|\nabla_z w_k|^2_{\mathrm{flat}}dz
\le C_{\mathrm{BG}}\int_M|\nabla_g w_k|^2dV_g
\le C'_{\mathrm{BG}}\|w_k\|_{H^1(M)}^2\le C''_{\mathrm{BG}},
\end{equation}
where $C_{\mathrm{BG}}$ is the bilipschitz constant of the Fermi chart
(Lemma~\ref{lem:atlas}, uniform under \textup{(BG$^{2+}$)}),
and the last two steps use the Poincar\'e inequality on $M$
and the graph-norm equivalence
$\|w_k\|_{H^1(M)}\lesssim \|w_k\|_{E,g}=1$
(Lemma~\ref{lem:graph-norm}(c)).
By the trace Sobolev inequality on the half-ball $B_{r_0/\varepsilon_k}^+$
(with the standard $H^1$ correction for bounded domains,
whose leading constant is independent of the radius by scaling),
\begin{equation}\label{eq:global-trace-bound}
\|\widetilde w_k(\cdot,0)\|_{L^q(\{|y'|<r_0/\varepsilon_k\})}
\le C_n\Big(\|\nabla\widetilde w_k\|_{L^2(B_{r_0/\varepsilon_k}^+)}
+\frac{\varepsilon_k}{r_0}\|\widetilde w_k\|_{L^2(B_{r_0/\varepsilon_k}^+)}\Big)
\le C,
\end{equation}
uniformly in $k$: the gradient term is $\le C_{\mathrm{BG}}^{1/2}$ by \eqref{eq:global-grad-bound},
and the $L^2$ correction satisfies
$(\varepsilon_k/r_0)\|\widetilde w_k\|_{L^2}
=r_0^{-1}\|w_k\|_{L^2(M)}\le r_0^{-1}\|w_k\|_{H^1}\le C'$
(by the critical scaling $\widetilde w_k(y)=\varepsilon_k^{(n-2)/2}w_k(\varepsilon_k y)$
and the graph-norm equivalence).
This global trace bound will be used in the tail estimate below.

By a standard diagonal argument over an exhaustion
$B_1^+\subset B_2^+\subset\cdots$,
we extract a subsequence and a function
$\widetilde w\in H^1_{\mathrm{loc}}(\R^n_+)$ such that
$\widetilde w_k\rightharpoonup\widetilde w$
weakly in $H^1(B_R^+)$ for every $R$.
By weak lower semicontinuity applied per ball and then
sending $R\to\infty$ (monotone convergence):
\begin{equation}\label{eq:wlimit-gradient-bound}
\|\nabla\widetilde w\|_{L^2(\R^n_+)}^2
:=\sup_R\|\nabla\widetilde w\|_{L^2(B_R^+)}^2
\le1,
\end{equation}
so $\widetilde w\in\dot H^1(\R^n_+)$.

We now show $\widetilde w\not\equiv0$.
Expanding the Hessian hypothesis \eqref{eq:contr-setup}:
\[
\big\langle D^2\mathcal J_g[U_{+,x_k,\varepsilon_k}]w_k,w_k\big\rangle
= \|w_k\|_{E,g}^2
   - (q-1)\lambda_k\int_{\partial M}U_{+,x_k,\varepsilon_k}^{q-2}w_k^2d\sigma_g
\to 0,
\]
where $\lambda_k:=\mathcal N^\circ_g(U_{+,x_k,\varepsilon_k})/\|U_{+,x_k,\varepsilon_k}\|_{L^q(\partial M)}^q\to S_\ast^{q/2}$
(the boundary Euler--Lagrange multiplier; note $\|U_+(\cdot,0)\|_{L^q}^q=S_\ast^{-q/2}$ by our normalization).
Since $\|w_k\|_{E,g}^2=1$, this forces
\begin{equation}\label{eq:bdry-mass}
(q-1)S_\ast^{q/2}\int_{\partial M}U_{+,x_k,\varepsilon_k}^{q-2}w_k^2d\sigma_g
\longrightarrow 1.
\end{equation}
Rescaling the boundary integral via the Fermi chart and
the near-isometry of the boundary Jacobian
(Lemma~\ref{lem:trace-no-linear}) gives
\[
\int_{\partial M}U_{+,x_k,\varepsilon_k}^{q-2}w_k^2d\sigma_g
=\int_{\R^{n-1}}U_+^{q-2}(y',0)
\widetilde w_k^2(y',0)dy'+o(1).
\]
(This is valid because $U_{+,x_k,\varepsilon_k}^{q-2}$
is concentrated at scale $\varepsilon_k$ near $x_k$,
so the contribution from outside the Fermi chart is negligible,
and inside the chart the boundary Jacobian correction is
$O(\varepsilon_k)$.)
Now, since $\widetilde w_k\rightharpoonup\widetilde w$ in
$H^1_{\mathrm{loc}}(\R^n_+)$ (from the diagonal extraction above), the traces satisfy
$\widetilde w_k(\cdot,0)\to\widetilde w(\cdot,0)$
strongly in $L^2_{\mathrm{loc}}(\R^{n-1})$ by the Rellich theorem.
Since $(n-2)(q-2)=2$, we have
$U_+^{q-2}(\cdot,0)\sim(1+|y'|)^{-2}\in L^{n-1}(\R^{n-1})$,
and $\|\widetilde w_k(\cdot,0)\|_{L^q}\le C_nC_{\mathrm{BG}}^{1/2}$
uniformly in $k$ by the global trace bound \eqref{eq:global-trace-bound}.
H\"older's inequality (with exponents $n-1$ and $q/2$,
since $\frac1{n-1}+\frac2q=1$) gives, for every $R>0$,
\[
\bigg|\int_{|y'|>R}U_+^{q-2}\widetilde w_k^2dy'\bigg|
\le \|U_+^{q-2}\|_{L^{n-1}(|y'|>R)}
     \|\widetilde w_k(\cdot,0)\|_{L^q}^2
\le C\|U_+^{q-2}\|_{L^{n-1}(|y'|>R)}\ \to\ 0
\]
as $R\to\infty$,
uniformly in $k$.
Combined with the $L^2_{\mathrm{loc}}$ convergence on $\{|y'|\le R\}$:
\[
\int_{\R^{n-1}}U_+^{q-2}\widetilde w_k^2dy'
\ \longrightarrow\
\int_{\R^{n-1}}U_+^{q-2}\widetilde w^2dy'.
\]
By \eqref{eq:bdry-mass}, the limit is
$\frac{1}{(q-1)S_\ast}>0$, so $\widetilde w\not\equiv0$.

\smallskip\noindent\textbf{Step 2 (Flat Hessian bound and spectral gap contradiction).}
By \eqref{eq:wlimit-gradient-bound} and the boundary integral convergence from Step~1,
\[
Q_\flat[\widetilde w]
:=\|\nabla\widetilde w\|_{L^2(\R^n_+)}^2
  -(q-1)S_\ast^{q/2}\int_{\R^{n-1}}U_+^{q-2}\widetilde w^2d\sigma
\ \le\ 1-1\ =\ 0.
\]
(Here $(q-1)S_\ast^{q/2}U_+^{q-2}=(q-1)S_\ast\hat U_+^{q-2}$
where $\hat U_+=S_\ast^{1/2}U_+$ is the $L^q$-normalized optimizer,
so the kernel and gap of $Q_\flat$ are those of
Lemma~\ref{lem:hessian-sphere} via the Cayley transfer.)
By Lemma~\ref{lem:hessian-sphere}
(which identifies the $n$-dimensional constrained kernel $E$
on the slice $T_{U_+}\mathcal S=\{\phi:\int U_+^{q-1}\phi=0\}$,
spanned by the dilation mode $\partial_\lambda U_+$ and the
tangential translation modes $\partial_{\xi'^\alpha}U_+$,
$\alpha=1,\dots,n-1$)
and Proposition~\ref{prop:robin-gap},
$Q_\flat\ge\gamma_\ast\|\cdot\|_{\dot H^1}^2$
on $T_{U_+}\mathcal S\cap E^\perp$.
The boundary constraint $\int U_+^{q-1}\widetilde w=0$
(from passage to the limit of the slice condition)
places $\widetilde w\in T_{U_+}\mathcal S$.
Since $\widetilde w\not\equiv0$ and
$Q_\flat[\widetilde w]\le0$,
$\widetilde w$ must have a nonzero component in the
constrained kernel $E$.
But the rescaled modulation orthogonality conditions
\eqref{eq:esc-orthog-gauge-final} pass to the limit
(weak convergence for the $\langle\cdot,\cdot\rangle_{\dot H^1}$
pairings with the decaying modes $\partial_\lambda U_+$,
$\partial_{\xi'}U_+$,
by the same Rellich+tail argument as above) and force
$\widetilde w\perp E$ in $T_{U_+}\mathcal S$.
Hence $\widetilde w\in T_{U_+}\mathcal S\cap E^\perp$, so
$Q_\flat[\widetilde w]\ge\gamma_\ast\|\widetilde w\|_{\dot H^1}^2>0$,
contradicting $Q_\flat[\widetilde w]\le0$.
\end{proof}

\end{document}
